%document
\documentclass[10pt]{beamer}
%theme
\usetheme{CambridgeUS}
% packages
\usepackage{color}
\usepackage{listings}
\usepackage[ngerman]{babel}
\usepackage[utf8]{inputenc}
\usepackage{multicol}


% color definitions
\definecolor{mygreen}{rgb}{0,0.6,0}
\definecolor{mygray}{rgb}{0.5,0.5,0.5}
\definecolor{mymauve}{rgb}{0.58,0,0.82}

\lstset{
    backgroundcolor=\color{white},
    % choose the background color;
    % you must add \usepackage{color} or \usepackage{xcolor}
    basicstyle=\footnotesize\ttfamily,
    % the size of the fonts that are used for the code
    breakatwhitespace=false,
    % sets if automatic breaks should only happen at whitespace
    breaklines=true,                 % sets automatic line breaking
    captionpos=b,                    % sets the caption-position to bottom
    commentstyle=\color{mygreen},    % comment style
    % deletekeywords={...},
    % if you want to delete keywords from the given language
    extendedchars=true,
    % lets you use non-ASCII characters;
    % for 8-bits encodings only, does not work with UTF-8
    frame=single,                    % adds a frame around the code
    keepspaces=true,
    % keeps spaces in text,
    % useful for keeping indentation of code
    % (possibly needs columns=flexible)
    keywordstyle=\color{blue},       % keyword style
    % morekeywords={*,...},
    % if you want to add more keywords to the set
    numbers=left,
    % where to put the line-numbers; possible values are (none, left, right)
    numbersep=5pt,
    % how far the line-numbers are from the code
    numberstyle=\tiny\color{mygray},
    % the style that is used for the line-numbers
    rulecolor=\color{black},
    % if not set, the frame-color may be changed on line-breaks
    % within not-black text (e.g. comments (green here))
    stepnumber=1,
    % the step between two line-numbers.
    % If it's 1, each line will be numbered
    stringstyle=\color{mymauve},     % string literal style
    tabsize=4,                       % sets default tabsize to 4 spaces
    % show the filename of files included with \lstinputlisting;
    % also try caption instead of title
    language = Java,
	showspaces = false,
	showtabs = false,
	showstringspaces = false,
	escapechar = ,
}

\def\ContinueLineNumber{\lstset{firstnumber=last}}
\def\StartLineAt#1{\lstset{firstnumber=#1}}
\let\numberLineAt\StartLineAt



\newcommand{\codeline}[1]{
	\alert{\texttt{#1}}
}

% This Document contains the information about this course.

% Authors of the slides
\author{Vincent Gerber, Tilman Hinnerichs}

% Name of the Course
\institute{Java Kurs}

% Fancy Logo
\titlegraphic{\hfill\includegraphics[height=1.25cm]{../templates/fsr_logo_cropped}}


\title{Java}
\subtitle{Controll Statements and Functions}
\date{\today}

\begin{document}

\begin{frame}
	\titlepage
\end{frame}
\begin{frame}{Overview}
	\setbeamertemplate{section in toc}[sections numbered]
	\tableofcontents
\end{frame}

\section{Recalling last session}
\begin{frame}{Conclusion}
	Datatypes
	\begin{itemize}
		\item int, long
		\item float, double
		\item String
	\end{itemize}
	Hello World example
	Switch variables example
\end{frame}

\section{Control Statements}
\begin{frame}{Control Statements}
	
	\begin{itemize}
		\item if, else, else if
		\item for
		\item while
	\end{itemize}
		
\end{frame}

\subsection{Ite}
\begin{frame}[fragile]{{\huge I}f {\huge T}hen {\huge E}lse}
\begin{lstlisting}
if(condition) {
	// do something if condition is true
} else if(another condition){
	// do if "else if" condition is true 
} else {
	// otherwise do this
}
\end{lstlisting}
\end{frame}

\begin{frame}[fragile]{{\huge I}f {\huge T}hen {\huge E}lse example}
\begin{lstlisting}
public class IteExample {

	public static void main(String[] args) {
		int myNumber = 5;
		
		if(myNumber == 3) {
			System.out.println("Strange number");
		} else if(myNumber == 2) {
			System.out.println("Unreachable code");
		} else {
			System.out.println("Will be printed");
		}
	}
    
}
\end{lstlisting}
\end{frame}

\begin{frame}{Conditions?}
How to compare things:
\begin{itemize}
    \item \texttt{==} Equal
    \item \texttt{!=} Not Equal
    \item \texttt{>} Greater Than
    \item \texttt{>=} Greater or Equal than
\end{itemize}

\textit{Note \#1}: Those statements return an boolean value \\
\textit{Note \#2}: You can concatenate multiple conditions\\ with \texttt{\&\&} (AND) or \texttt{||} (OR)
\end{frame}

\subsection{for}
\begin{frame}[fragile]{for}
\begin{lstlisting}
for(initial value, condition, change) {
	// do code while condition is true
}
\end{lstlisting}
\end{frame}

\begin{frame}[fragile]{for example}
\begin{lstlisting}
public class ForExample {

	public static void main(String[] args) {
		for(int i = 0; i <= 10; i++) {
			System.out.print("na ");
		}
		System.out.println("BATMAN!");
	}
    
}
\end{lstlisting}
\end{frame}

\begin{frame}[fragile]{Another for example}
	\begin{lstlisting}
public class AnotherForExample {

	public static void main(String[] args) {
		for(int i = 0; i <= 10; i++) {
			System.out.println(i);
		}
	}
}	
	\end{lstlisting}
\end{frame}

\subsection{while}
\begin{frame}[fragile]{while}
\begin{lstlisting}
while(condition) {
	// do code while condition is true
}
\end{lstlisting}
\end{frame}
	
\begin{frame}[fragile]{while example}
\begin{lstlisting}
public class WhileExample {

	public static void main(String[] args) {
		int a = 0;
		while(a <= 10) {
			System.out.println(a);
            a++; // Otherwise you would get an endless loop
		}
	}
    
}
\end{lstlisting}
\end{frame}

\begin{frame}[fragile]{Functions}
	\begin{lstlisting}
Type func_name (type parameters, ...) {
	function statements;
	...
	return value;
}
	\end{lstlisting}
\end{frame}
\begin{frame}[fragile]{Functions example}
	\begin{lstlisting}
int add(int a, int b) {
	int c = a + b;
	System.out.println(c);
	return c;
}
	\end{lstlisting}
\end{frame}

\begin{frame}[fragile]{Functions example}
	\begin{lstlisting}
int add(int a, int b) {
	return a+b;
}
	\end{lstlisting}
\end{frame}

\begin{frame}[fragile]{Well known functions example}
	You all have written that before:
	\begin{lstlisting}
	public static void main(String[] args) {
		...
	}
	\end{lstlisting}
\end{frame}

\begin{frame}[fragile]{Exercises}
	\textbf{1. Display the n terms of odd natural number.}
	\begin{example}	 	
	Input: 10, Output: 1, 3, 5, 7, 9
	\end{example}
	
	\textbf{2. Find the sum and average of 3 numbers.}
	\begin{example}	
	Input: 6, 4, 3 Output: 13 and 4,3
	\end{example}	
	
	\textbf{3. Find the smallest number among three numbers.}
	\begin{example}	
	Input: 6, 4, 3 Output: 3 
	\end{example}
	
\end{frame}

\begin{frame}[fragile]{Adv. Exercises}
	\textbf{4. Display the pattern like a diamond.}
	\begin{example}	
	\begin{verbatim}	
	      *                                                                          
     ***                                                                         
    *****                                                                        
   *******                                                                       
  *********                                                                      
 ***********                                                                     
*************                                                                    
 ***********                                                                     
  *********                                                                      
   *******                                                                       
    *****                                                                        
     ***                                                                         
      * 
	\end{verbatim}
	\end{example}
\end{frame}

\begin{frame}[fragile]{Adv. Exercise}
	\textbf{5. Find all primes up to n.}
	\begin{example}	
	Input: 11 Output: 2, 3, 5, 7, 11
	\end{example}
\end{frame}

\ifx true false
\section{OOP in Java}

\begin{frame}{}
	\begin{center}
		{\huge Object Oriented Programming}
	\end{center}
\end{frame}

\subsection{General information}

\begin{frame}[fragile]{Class Student}
\begin{lstlisting}
public class Student {

	// Attributes
	private String name; 
	private int matriculationNumber; 
	
	
	// Methods
	public void setName(String name) {
		this.name = name;
	}
    
	public int getMatriculationNumber() {
		return matriculationNumber;
	}
    
}
\end{lstlisting}

% What is visible here:
% Attributes store the state of the object
% Methods implement the behaviour of the object

\end{frame}

\begin{frame}[fragile]{Creation}
	We learned how to declare and assign a primitive datatype.}
	
	\begin{lstlisting}
	    int a; // declare a
	    a = 273; // assign 273 to a
	\end{lstlisting} 
	
	The creation of an object works similar.
	
	\begin{lstlisting}
	    Student example = new Student(); 
		// create an instance of Student
	\end{lstlisting}
	The \textbf{object} derived from a \textbf{class} is also called \textbf{instance}.
	The variable is called the \textbf{reference}.
\end{frame}

\subsection{Methods}
\begin{frame}[fragile]{Calling a Method}
	\begin{lstlisting}
	public class Student {
    
		private String name;
	
		public String getName() {
        	return name;
    	}
		
	    public void setName(String newName) {
			name = newName;
	    }
	   
	}
	\end{lstlisting}
	The class \emph{Student} has two methods: \emph{void printTimetable()} and \emph{void printName()}.
\end{frame}

\begin{frame}[fragile]{Calling a Method}
\begin{lstlisting}
public class Main {

    public static void main(String[] args) {
        Student example = new Student(); // creation
        example.setName("Jane"); // method call
        String name = example.getName(); 
		System.out.println(name); // Prints "Jane"
    }
    
}
	\end{lstlisting}
	You can call a method of an object after its creation with \textbf{reference.methodName();}.
\end{frame}

\begin{frame}[fragile]{Calling a Method}
\begin{lstlisting}
	public class Student {
    
		private String name;
	
	    public void setName(String newName) {
			name = newName;
			printName();   // Call own method
			this.printName(); // Or this way
	    }
	    
	    public void printName() {
	        System.out.println(name);
	    }
        
	}
	\end{lstlisting}
	You can call a method of the own object by simply writing \textbf{methodName();} or \textbf{this.methodName();}
\end{frame}

\begin{frame}[fragile]{Methods with Arguments}

\begin{lstlisting}
public class Calc {

    public void add(int summand1, int summand2) {
        System.out.println(summand1 + summand2);
    }
	    
    public static void main(String[] args) {
        int summandA = 1;
        int summandB = 2;
        Calc calculator = new Calc();
        System.out.print("1 + 2 = ");
        calculator.add(summandA, summandB); 
		// prints: 3
    }
    
}
	\end{lstlisting}
\end{frame}

\subsection{Return Value}
\begin{frame}[fragile]{Methods with Return Value}
	A method without a return value is indicated by \textbf{void}:
	\begin{lstlisting}
	public void add(int summand1, int summand2) {
	    System.out.println(summand1 + summand2);
	}
	\end{lstlisting}
	A method with an \textbf{int} as return value:
	\begin{lstlisting}
	public int add(int summand1, int summand2) {
	    return summand1 + summand2;
	}
	\end{lstlisting}
	% TODO explain return statement
\end{frame}

\begin{frame}[fragile]{Calling Methods with a return value}
	\begin{lstlisting}
	public class Calc {
	
	    public int add(int summand1, int summand2) {
	        return summand1 + summand2;
	    }
	    
	    public static void main(String[] args) {
	        Calc calculator = new Calc();
	        int sum = calculator.add(3, 8);
	        System.out.print("3 + 8 = " + sum); 
			// prints: 3 + 8 = 11
	    }
        
	}
	\end{lstlisting}
\end{frame}

\subsection{Constructor}

\begin{frame}[fragile]{Constructors}
	\begin{lstlisting}
	public class Calc {
    
		private int summand1;
		private int summand2;
	
	    public Calc() {
			summand1 = 0;
			summand2 = 0;
	    }
        
	}
	\end{lstlisting}
	A constructor gets called upon creation of the object
\end{frame}

\begin{frame}[fragile]{Constructors with Arguments}
	\begin{lstlisting}
	public class Calc {
    
		private int summand1;
		private int summand2;
	
	    public Calc(int x, int y) {
			summand1 = x;
			summand2 = y;
	    }
        
	}
	\end{lstlisting}
	\begin{lstlisting}
	[...]
	Calc myCalc = new Calc(7, 9);
	\end{lstlisting}
    
	A constructor can have arguments as well!
\end{frame}

\section{Conclusion}
\subsection{An Example}

\begin{frame}{An Example}
	You want to program an enrollment system, for a programming course. \\
	\vspace{1em}
	Your classes are:\\
	\begin{description}
		\item[student] who wants to attend the course
		\item[lesson] which is a part of the course
		\item[tutor] the guy with the bandshirt
		\item[room] where your lessons take place
		\item[\dots]
	\end{description}
% 	The more you think about it, the more complex this program becomes.
% 	Focus on the relevant things.
%	Think about how the objects can be in relation, this will be discussed later
%	Show prepared classes in Java
\end{frame}


\begin{frame}[fragile]{Class Student}
\begin{lstlisting}
public static void main(String[] args) {
	Student peter = new Student();
	peter.changeName("Peter");
}
\end{lstlisting}
\end{frame}
\fi

\end{document}